\scanpage{35}
(\uline{5 bis}) \uline{Détermination de} $\Sigma^n(F)$ (2\textsuperscript{ème} mouture \ldots)

Considérons l'extension vectorielle définissant $F$ \inserted{un fibré affine de rg 1}
\[
0\longrightarrow L\xrightarrow{i_1}\deleted{\unsure{B}}\overset{\deleted{\unsure{6|e|\leq}}}{V}\xrightarrow{\pi_1}\mathcal O_S\longrightarrow0 \qquad (V=\check B) \tag*{(47)}
\]
\[
 F=\pi_1^{-1}(1)\subset\check{\mathbb V}(V) \tag*{(48)}
\]
on en déduit, pour tout $n\geq1$, une structure
d'extension
\[
0\longrightarrow L\otimes\underline{\mathrm{Sym}}^{n-1}(\deleted{\unsure{B}}V)\xrightarrow{i_n}\underline{\mathrm{Sym}}^n(\deleted{\unsure{\check B}}V)\xrightarrow{\pi_n}\mathcal O_S\longrightarrow0 \tag*{(49)}
\]
qui \inserted{(en \unsure{itérant|théorie})} met sur $\mathrm{Sym}^n(V)$ une filtration
\deleted{\unsure{dont les}} décroissante par les
\[
\underbrace{L^i\otimes\underline{\mathrm{Sym}}^{n-i}(V)}_{\mathrm{Fil}^i\,\mathrm{Sym}^n(V)}\subset\underline{\mathrm{Sym}}^n(V)\quad(0\leq i\leq n) \tag*{(50)}
\]

\deleted{Considérons}
\[
\deleted{\Sigma^n\deleted{\unsure{(F)|F}}=\pi_n^{-1}(\deleted{\unsure{\check{\mathbb V}(\mathcal O_S)}}1)\subset\check{\mathbb V}(\underline{\mathrm{Sym}}^n(V))} \tag*{\deleted{(51)}}
\]
\deleted{le fibré affine associé à l'extension (50).}
\deleted{On va définir un isomorphisme canonique}
\[
\deleted{\underbrace{\Sigma^n(F)}_{\substack{n.\text{ième puiss.}\\\text{symétrique}}}\xrightarrow[\sim]{\varphi_n}\Sigma^n} \tag*{\deleted{(52)}}
\]

\deleted{\unsure{1)|6)} Description ``fonctorielle''. On définira le
morphisme $\varphi_n$ par passage au quotient du}
\[
\deleted{\left\{\begin{aligned}
 \psi_n &: F^n\longrightarrow\Sigma^n\\
 \psi_n(x_1,\ldots,x_n)&=x_1
\end{aligned}\right.} \tag*{\deleted{(53)}}
\]
% archive sheet number: 34
